\documentclass[11pt]{amsart}

\usepackage[utf8]{inputenc}
\usepackage{amssymb}
\usepackage{mathtools}
\usepackage[expansion=false]{microtype}
\usepackage[shortlabels]{enumitem}
\usepackage{xcolor}
\usepackage[
  pdftitle={Ap\'ery polynomials as Stieltjes moment sequences},
  pdfauthor={Knowledge Sense},
  pdfsubject={Mathematics},
  pdfkeywords={Apery polynomials, Stieltjes moment problem, Stieltjes transform, complete elliptic integral, lattice Green function},
  colorlinks=true, linkcolor=blue!50!black, citecolor=blue!50!black, urlcolor=blue!50!black
]{hyperref}

\hyphenation{manuscript}
\emergencystretch=2em

\numberwithin{equation}{section}

\newtheorem{theorem}{Theorem}[section]
\newtheorem{proposition}[theorem]{Proposition}
\newtheorem{lemma}[theorem]{Lemma}
\newtheorem{corollary}[theorem]{Corollary}
\theoremstyle{definition}
\newtheorem{definition}[theorem]{Definition}
\theoremstyle{remark}
\newtheorem{remark}[theorem]{Remark}

\newcommand{\C}{\mathbb{C}}
\newcommand{\R}{\mathbb{R}}
\newcommand{\E}{\mathbb{E}}
\newcommand{\Pim}{\Pi_{-}}
\newcommand{\Pip}{\Pi_{+}}
\newcommand{\dd}{\,\mathrm{d}}
\DeclareMathOperator{\re}{Re}
\DeclareMathOperator{\im}{Im}
\newcommand{\Kc}{\mathbf{K}}

\begin{document}

\title[Ap\'ery polynomials as Stieltjes moment sequences]{Ap\'ery polynomials as Stieltjes moment sequences}

\author{Knowledge Sense}
\address{Knowledge Sense}
\email{math@knowledgesense.jp}

\date{13 October 2026}

\subjclass[2020]{Primary 44A60; Secondary 05A10, 11B65, 30E05, 33E05}

\keywords{Ap\'ery polynomials, Stieltjes moment sequence, Stieltjes transform, analytic continuation, complete elliptic integral}

\begin{abstract}
For real $x$ let $A_n(x)=\sum_{j=0}^{n}\binom{n}{j}^2\binom{n+j}{j}^2x^j$ be the
Ap\'ery polynomials. Sokal conjectured that $(A_n(x))_{n\ge0}$ is a Stieltjes
moment sequence for every $x\ge1$; the case $x=1$, concerning the Ap\'ery
numbers, was proved by Edgar. We prove the conjecture. For every real $x\ge1$
there is a probability measure $\mu_x$ on $[0,c_x]$, where
$c_x=\bigl(1+2\sqrt{x}+2\sqrt{x+\sqrt{x}}\bigr)^2$, such that
$A_n(x)=\int t^n\dd\mu_x(t)$ for all $n\ge0$. The measure is absolutely
continuous, and its density is given explicitly by one-dimensional integrals
of complete elliptic integrals; it is strictly positive on $(0,c_x)$ and has
at most a squared logarithmic singularity at the origin. The proof writes the
generating function of $x^{-n}A_n(x)$ as an average of a resolvent over a
three-dimensional torus, continues the resulting Stieltjes transform through a
Joukowski-type change of variables, and identifies the boundary values on the
second sheet by means of the logarithmic monodromy of the complete elliptic
integral.
\end{abstract}

\maketitle

\section{Introduction}

A sequence $(a_n)_{n\ge0}$ of real numbers is a \emph{Stieltjes moment
sequence} if there is a positive Borel measure $\mu$ on $[0,\infty)$ such that
$a_n=\int t^n\dd\mu(t)$ for all $n\ge0$; equivalently, the Hankel matrix
$(a_{i+j})_{i,j\ge0}$ is totally positive \cite{WZ}. Many sequences
arising in enumerative combinatorics are Stieltjes moment sequences; see \cite{So20,WZ}
and the references there.

The Ap\'ery numbers $A_n=\sum_{k=0}^{n}\binom{n}{k}^2\binom{n+k}{k}^2$ appear in
Ap\'ery's proof of the irrationality of $\zeta(3)$ \cite{Ap}. Sokal conjectured
that they form a Stieltjes moment sequence, and Edgar \cite{Ed} proved this,
with a representing density expressed through Heun functions. The
\emph{Ap\'ery polynomials}
\begin{equation}\label{eq:apery}
A_n(x)=\sum_{j=0}^{n}\binom{n}{j}^2\binom{n+j}{j}^2x^j\qquad(n\ge0)
\end{equation}
satisfy $A_n(1)=A_n$. Sokal \cite[Introduction, (1.2)]{So20} recorded as open
the more general conjecture that $(A_n(x))_{n\ge0}$ is a Stieltjes moment
sequence for every $x\ge1$, and the conjecture is restated as Conjecture~1 in
\cite[slide~29]{So25}. Our main result settles it.

\begin{theorem}\label{thm:main}
Let $x\ge1$ be real and put
\begin{equation}\label{eq:cx}
c_x=\Bigl(1+2\sqrt{x}+2\sqrt{x+\sqrt{x}}\Bigr)^2 .
\end{equation}
There is a unique positive Borel measure $\mu_x$ on $\R$ such that
\[
A_n(x)=\int t^n\dd\mu_x(t)\qquad\text{for all } n\ge0 .
\]
It is a probability measure, it is absolutely continuous with respect to
Lebesgue measure, and its support is $[0,c_x]$. Its density $\varphi_x$,
given explicitly in \eqref{eq:density} below, is strictly positive on
$(0,c_x)$ and satisfies $\varphi_x(t)=O(\log^2t)$ as $t\to0^+$. In particular,
$(A_n(x))_{n\ge0}$ is a Stieltjes moment sequence.
\end{theorem}

\subsection*{The density}
We use the complete elliptic integral of the first kind in the
\emph{parameter} convention,
\begin{equation}\label{eq:K}
\Kc(m)=\int_0^{\pi/2}\frac{\dd\theta}{\sqrt{1-m\sin^2\theta}}\qquad(m<1),
\end{equation}
so that $\Kc(m)$ is $K(\sqrt m)$ in the modulus notation of \cite{DLMF}. Let
$x\ge1$ and put $b=2x^{-1/2}\in(0,2]$ and $B=4+b$. Define
\begin{equation}\label{eq:f}
f(c)=\frac{1}{2\pi^2}\,\Kc\Bigl(1-\frac{c^2}{16}\Bigr)\quad(0<|c|<4),\qquad
f(c)=0\quad(|c|>4),
\end{equation}
and, for real $w$,
\begin{equation}\label{eq:h}
h(w)=\frac1\pi\int_{-b}^{b}\frac{f(w-u)}{\sqrt{b^2-u^2}}\dd u .
\end{equation}
For real $w>b$ define further
\begin{align}
T(w)&=\frac{1}{2\pi^3}\int_{-b}^{b}\Kc\Bigl(1-\frac{(w-u)^2}{16}\Bigr)
\frac{\dd u}{\sqrt{b^2-u^2}},\label{eq:T}\\
J(w)&=\frac{1}{\pi^3}\int_{b}^{w}\kappa\Bigl(\frac{(w-u)^2}{16}\Bigr)
\frac{\dd u}{\sqrt{u^2-b^2}},\label{eq:J}
\end{align}
where
\begin{equation}\label{eq:kappa}
\kappa(m)=\Kc(m)\quad(0\le m<1),\qquad
\kappa(m)=m^{-1/2}\,\Kc(1/m)\quad(m>1).
\end{equation}
Finally let $r_x(t)=x^{-1/2}\bigl(\sqrt t+1/\sqrt t\bigr)$ and
\begin{equation}\label{eq:density}
\varphi_x(t)=
\begin{cases}
\dfrac{h(r_x(t))}{\sqrt{xt}}, & 1\le t<c_x,\\[2ex]
\dfrac{2T(r_x(t))-h(r_x(t))+2J(r_x(t))}{\sqrt{xt}}, & 0<t<1,
\end{cases}
\end{equation}
and $\varphi_x(t)=0$ for $t\notin(0,c_x)$. The values of $f$ at the finitely
many excluded points are irrelevant. All integrals converge absolutely: the
only singularities are inverse square roots at the endpoints and logarithmic
singularities of $\Kc$ at the parameter value $1$. The function $r_x$
decreases from $+\infty$ to $b$ on $(0,1]$ and increases from $b$ to $B$ on
$[1,c_x]$; at $t=1$ both lines of \eqref{eq:density} equal $h(b)/\sqrt x$.

The positivity of $\varphi_x$ is elementary. Since $f>0$ on $(-4,4)$, we
have $h>0$ on $(-B,B)$. For $w>b$ the integrand of $T(w)$ is positive on all
of $(-b,b)$, whereas $h(w)$ retains it only where $w-u<4$; thus $0\le h\le
T$, and $\kappa>0$ gives $J\ge0$. Hence $\varphi_x(t)\ge
T(r_x(t))/\sqrt{xt}>0$ for $0<t<1$. What has to be proved is that the
measure $\varphi_x(t)\dd t$ has the moments $A_n(x)$, with exactly the
combination $2T-h+2J$ on $(0,1)$.

\subsection*{Relation to earlier work}
For $x=1$ Theorem~\ref{thm:main} is Edgar's theorem \cite[Theorem~1]{Ed},
and $c_1=17+12\sqrt2$ is the endpoint found there \cite[Notation~3]{Ed}; by
uniqueness, $\mu_1$ is Edgar's measure. The present paper extends the
existence of a representing measure to every $x>1$ and gives the explicit
density \eqref{eq:density}; the argument treats all $x\ge1$ in the same way.

The proof starts from the polynomials $D_n(\alpha)=\alpha^nA_n(1/\alpha)$.
By Straub's theorem on multivariate Ap\'ery numbers \cite[Theorem~3.1]{St},
applied with the partition $\lambda=(2,2)$, $D_n(\alpha)$ is the diagonal
coefficient of $u^nv^nw^nz^n$ in
\[
\frac{1}{(1-u-v)(1-w-z)-\alpha uvwz}.
\]
An average of this rational function over a torus leads to the spectral
representation of Proposition~\ref{prop:spectral} (Remark~\ref{rem:torus}).

Neuschel \cite[Theorem~4.1]{Ne} determined the asymptotic behaviour of the
Ap\'ery polynomials in $\C\setminus(-\infty,0]$; in particular
$A_n(x)^{1/n}\to c_x$ for $x>0$. For a positive measure with compact support
the $n$th roots of the moments converge to the largest point of the support,
so once a representing measure is known to exist, the right endpoint $c_x$ of
its support is already determined by \cite{Ne}.

Two related questions are not addressed here. Sokal also conjectures that
$(A_n(x))_{n\ge0}$ is not a Stieltjes moment sequence for $0<x<1$
\cite[slide~29]{So25}; in our proof the hypothesis $x\ge1$ is essential for
the holomorphic continuation of the spectral density across the central
interval (Remark~\ref{rem:xge1}). His Conjecture~2 in \cite[slide~29]{So25}
asserts that $(A_n(1+y))_{n\ge0}$ is coefficientwise Hankel-totally positive
in $y$, that is, that the minors of its Hankel matrix are polynomials in $y$
with nonnegative coefficients. This would imply that $(A_n(x))_{n\ge0}$ is a
Stieltjes moment sequence for every $x\ge1$, but it is a different,
coefficientwise statement, and it remains open.

\subsection*{Strategy}
Write $\alpha=1/x\in(0,1]$, so that $b=2\sqrt\alpha$, and let $D_n(\alpha)=\alpha^nA_n(1/\alpha)$.
\begin{enumerate}[(i)]
\item \emph{A spectral representation.} Let $L=4\cos\theta_1\cos\theta_2+b\cos\theta_3$ with independent
uniform angles. Then
\[
\textstyle\sum_n D_n(\alpha)\tau^{2n}=\E\bigl[(1-\tau L+\alpha\tau^2)^{-1}\bigr]
\]
(Proposition~\ref{prop:spectral}), and $L$ has the density $h$.
\item \emph{A folded Stieltjes transform.} With $s=\tau^{-2}$ this becomes
\[
\textstyle\sum_n D_n(\alpha)s^{-n-1}=g(r)/\sqrt s,
\]
where $g$ is the Cauchy transform of
$h$ and $r=\sqrt s+\alpha/\sqrt s$. The map $s\mapsto r$ is two-to-one: it
sends the upper half-disc $\{|s|<\alpha,\ \im s>0\}$ into the lower half-plane
and the semicircle $|s|=\alpha$ to the interval $(0,b)$. Continuing the
transform into the half-disc therefore requires continuing $g$ across $(0,b)$,
that is, a holomorphic continuation of $h$ from $(-b,b)$.
\item \emph{The continuation.} In Section~\ref{sec:continuation} we construct
this continuation $H$ explicitly in the lower half-plane. On the central
interval $(-b,b)$ the condition $b\le2$ keeps all differences of points of
$(-b,b)$ inside $(-4,4)$, and the jump of $\Kc$ across its branch cut, with
its exact sign, shows that $H$ is real there. To the right of $b$ the
boundary values of $H$ have real part $T+J$.
\item \emph{Reconstruction.} In Section~\ref{sec:proof} the Stieltjes transform
is shown to be holomorphic on $\C\setminus[0,S]$, with growth bounds at the
origin that exclude atoms, and Cauchy's integral formula recovers the measure
from the jump across the cut.
\end{enumerate}

\subsection*{Dependencies and organization}
Besides elementary complex analysis, the proof uses only classical facts on
the complete elliptic integral and the hypergeometric equation, which we
collect in Section~\ref{sec:elliptic} with references to \cite{DLMF}. The
results of \cite{Ed,Ne,St} and of \cite{DJ01,DJ06} recalled in this paper
are not used in the proofs: Proposition~\ref{prop:spectral} is verified
directly from the coefficients, and the endpoint $c_x$ is obtained from the
construction. Section~\ref{sec:spectral} proves the spectral
representation and the regularity of $h$. Section~\ref{sec:continuation}
constructs the continuation $H$ and proves the bounds needed later.
Section~\ref{sec:proof} studies the Stieltjes transform and proves
Theorem~\ref{thm:main}. Section~\ref{sec:remarks} contains further remarks.

\section{The complete elliptic integral}\label{sec:elliptic}

For $m\in\C\setminus[1,\infty)$ and $\theta\in[0,\pi/2]$ the number
$1-m\sin^2\theta$ does not lie in $(-\infty,0]$. With the principal square
root, \eqref{eq:K} therefore defines a holomorphic function on
$\C\setminus[1,\infty)$, which is positive on $(-\infty,1)$. Expanding the
integrand gives $\Kc(m)=\frac\pi2\sum_{j\ge0}\bigl(\frac{(1/2)_j}{j!}\bigr)^2m^j$
for $|m|<1$, so $\frac2\pi\Kc$ is the Gauss function
${}_2F_1(\frac12,\frac12;1;m)$ and $\Kc$ satisfies the hypergeometric equation
\begin{equation}\label{eq:hyp}
m(1-m)y''+(1-2m)y'-\tfrac14y=0
\end{equation}
\cite[Section~15.10]{DLMF}. Its singular points $0$, $1$, $\infty$ are regular, with
exponents $\{0,0\}$, $\{0,0\}$ and $\{\frac12,\frac12\}$ respectively. In
particular $\Kc$ continues analytically across $(1,\infty)$ from either side,
so the boundary values $\Kc(m\pm i0)=\lim_{\varepsilon\downarrow0}\Kc(m\pm i\varepsilon)$
exist and are continuous for $m>1$.

\begin{lemma}[Logarithmic expansion]\label{lem:log}
There is a power series $E(\varepsilon)=\sum_{j\ge0}e_j\varepsilon^j$ with real
coefficients, convergent for $|\varepsilon|<1$, such that
\begin{equation}\label{eq:logexp}
\Kc(m)=\frac1\pi\,\Kc(1-m)\log\frac{1}{1-m}+E(1-m)
\end{equation}
for $|1-m|<1$, $m\notin[1,\infty)$, with the principal logarithm.
\end{lemma}

\begin{proof}
In the modulus notation, \cite[19.12.1]{DLMF} states that
\[
K(k)=\sum_{j\ge0}\Bigl(\frac{(1/2)_j}{j!}\Bigr)^2k'^{2j}\Bigl(\log\frac1{k'}+d(j)\Bigr)
\]
for $0<k<1$, with $k'=\sqrt{1-k^2}$ and explicit real numbers $d(j)$, the
series converging for $|k'|<1$. Put $m=k^2$, so that $k'^2=1-m$ and
$\log(1/k')=\frac12\log\frac1{1-m}$. Since
$\sum_j\bigl(\frac{(1/2)_j}{j!}\bigr)^2(1-m)^j=\frac2\pi\Kc(1-m)$, this is
\eqref{eq:logexp} for $0<m<1$, with
$e_j=\bigl(\frac{(1/2)_j}{j!}\bigr)^2d(j)$. Both sides are holomorphic on
$\{|1-m|<1\}\setminus[1,2)$, so the identity persists there.
\end{proof}

\begin{lemma}[Real part on the cut]\label{lem:cut}
For $m>1$,
\[
\re\Kc(m+i0)=\re\Kc(m-i0)=m^{-1/2}\,\Kc(1/m)>0 .
\]
\end{lemma}

\begin{proof}
Let $\theta_0=\arcsin(m^{-1/2})$. As $\varepsilon\downarrow0$, the
integrand $(1-(m\pm i\varepsilon)\sin^2\theta)^{-1/2}$ converges to the positive number
$(1-m\sin^2\theta)^{-1/2}$ for $\theta<\theta_0$, and to a purely imaginary
number for $\theta>\theta_0$. Since
$|1-(m\pm i\varepsilon)\sin^2\theta|\ge|1-m\sin^2\theta|$, the integrands are
dominated by the integrable function $|1-m\sin^2\theta|^{-1/2}$. Hence, with
the substitution $\sin\theta=m^{-1/2}\sin\phi$,
\[
\re\Kc(m\pm i0)=\int_0^{\theta_0}\frac{\dd\theta}{\sqrt{1-m\sin^2\theta}}
=\frac{1}{\sqrt m}\int_0^{\pi/2}\frac{\dd\phi}{\sqrt{1-m^{-1}\sin^2\phi}} .
\qedhere
\]
\end{proof}

The full boundary values are $\Kc(m\pm i0)=m^{-1/2}\bigl(\Kc(1/m)\pm
i\,\Kc(1-1/m)\bigr)$, a form of the reciprocal-modulus transformation
\cite[19.7.3]{DLMF}; only the real part will be needed.

\begin{lemma}[Bounds]\label{lem:bounds}
There is an absolute constant $C$ such that, for $m\in\C\setminus[1,\infty)$
and also for the boundary values $\Kc(m\pm i0)$ with $m>1$,
\[
|\Kc(m)|\le
\begin{cases}
C\bigl(1+\bigl|\log|1-m|\bigr|\bigr), & |m|\le4,\\[0.5ex]
C\,(1+\log|m|)\,|m|^{-1/2}, & |m|\ge4.
\end{cases}
\]
\end{lemma}

\begin{proof}
For $|1-m|\le\frac12$ the first bound follows from \eqref{eq:logexp}, since
$\Kc(1-m)$ and $E(1-m)$ are bounded there and
$|\log\frac1{1-m}|\le|\log|1-m||+\pi$. On the remaining part of
$\{|m|\le4\}$, $\Kc$ is continuous up to both sides of the cut, hence bounded.
At $\infty$ the exponents of \eqref{eq:hyp} are $\frac12,\frac12$, and the
nearest other singular point is $1$. Thus on $\{|m|>2\}\setminus[2,\infty)$
every solution of \eqref{eq:hyp} is a linear combination of
$(-m)^{-1/2}\Phi_1(1/m)$ and
$(-m)^{-1/2}\bigl(\Phi_1(1/m)\log(-m)+\Phi_2(1/m)\bigr)$, where $\Phi_1,\Phi_2$
are holomorphic in $|\zeta|<1$ and the powers and the logarithm are principal.
Applied to $\Kc$, this gives the second bound, including the boundary values on
both sides of $[2,\infty)$.
\end{proof}

We next continue $m\mapsto\Kc(m)$ along $m=1-z^2/16$. Let
$\Pim=\{z:\im z<0\}$ and $\Pip=\{z:\im z>0\}$. For $\re z>0$ and $\im z<0$ we
have $z^2\notin(-\infty,0]$, hence $1-z^2/16\notin[1,\infty)$.

\begin{lemma}[The continued kernel]\label{lem:P}
There is a unique holomorphic function $P$ on $\Pim$ such that
$P(z)=\frac{1}{2\pi^2}\Kc(1-z^2/16)$ when $\re z>0$. It has the following
properties.
\begin{enumerate}[(a)]
\item $P$ extends continuously to $\overline{\Pim}\setminus\{0,-4\}$, and for
$c>0$ its boundary value is $P(c)=\frac1{2\pi^2}\Kc(1-c^2/16)>0$.
\item For $0<c<4$,
\begin{equation}\label{eq:jump}
P(-c-i0)=f(c)+\frac{i}{\pi^2}\,\Kc\Bigl(\frac{c^2}{16}\Bigr).
\end{equation}
\item There is a constant $C$ such that for $z\in\overline{\Pim}\setminus\{0,-4\}$
\[
|P(z)|\le
\begin{cases}
C\bigl(1+\bigl|\log|z|\bigr|+\bigl|\log|z+4|\bigr|\bigr), & |z|\le8,\\[0.5ex]
C\,(1+\log|z|)\,|z|^{-1}, & |z|\ge8.
\end{cases}
\]
\end{enumerate}
\end{lemma}

\begin{proof}
Substituting $m=1-z^2/16$ into \eqref{eq:hyp} gives a linear differential
equation of second order in $z$ whose coefficients are holomorphic outside
$\{0,4,-4\}$. The function $z\mapsto\Kc(1-z^2/16)$ solves it on
$\{\re z>0,\ \im z<0\}$. Since $\Pim$ is simply connected and contains no
singular point, this solution continues uniquely to a holomorphic function on
$\Pim$, and it continues analytically across every real point other than $0$
and $\pm4$. Near $z=4$ the continuation is the principal function
$\Kc(1-z^2/16)$, which is holomorphic at $m=0$. This gives the continuity in
(a). For $c>0$ we have $1-c^2/16<1$, where $\Kc$ is holomorphic and positive.

For (b), let $z\in\Pim$, $|z|<4$ and $\re z>0$. Then $1-m=z^2/16$, and
$\log\frac1{1-m}=2\log\frac4z$ with principal logarithms, because
$\arg z\in(-\pi/2,0)$. Lemma~\ref{lem:log} gives
\begin{equation}\label{eq:Plocal}
2\pi^2P(z)=\frac2\pi\,\Kc\Bigl(\frac{z^2}{16}\Bigr)\log\frac4z+E\Bigl(\frac{z^2}{16}\Bigr).
\end{equation}
With the principal branch of $\log z$, the right-hand side is holomorphic on
the whole half-disc $\Pim\cap\{|z|<4\}$, so \eqref{eq:Plocal} holds there. As
$z\to c\in(0,4)$, \eqref{eq:Plocal} reads
$\Kc(1-c^2/16)=\frac2\pi\Kc(c^2/16)\log\frac4c+E(c^2/16)$. As $z\to-c$ from
$\Pim$, we have $\log z\to\log c-i\pi$, so $\log\frac4z\to\log\frac4c+i\pi$,
while $\Kc(z^2/16)\to\Kc(c^2/16)$ and $E(z^2/16)\to E(c^2/16)$. Subtracting the
two limits gives
$2\pi^2P(-c-i0)=\Kc(1-c^2/16)+2i\,\Kc(c^2/16)$, which is \eqref{eq:jump}.

For (c), \eqref{eq:Plocal} gives the bound near $z=0$. Near $z=-4$ we have
$m\to0$, a regular singular point with exponents $\{0,0\}$, so every
solution is $O(1+|\log|m||)$ there, and $|m|\asymp|z+4|$. Elsewhere on
$\overline{\Pim}\cap\{|z|\le8\}$, $P$ is continuous. For $|z|\ge8$ we have
$|m|\ge3$. The set $\overline{\Pim}\cap\{|z|\ge8\}$ is simply connected, and on
it $P$ is a fixed linear combination of the two solutions at $\infty$ used in
the proof of Lemma~\ref{lem:bounds}, with $(-m)^{-1/2}$ and $\log(-m)$
continued along this set. Since $\arg(-m)=\arg(z^2/16-1)$ varies in a bounded
interval there, we obtain $|P(z)|\le C(1+\log|m|)|m|^{-1/2}\le
C'(1+\log|z|)/|z|$.
\end{proof}

\section{The spectral representation}\label{sec:spectral}

Throughout this section $0<\alpha\le1$, $b=2\sqrt\alpha$ and $B=4+b$. Let
\begin{equation}\label{eq:D}
D_n(\alpha)=\alpha^nA_n(1/\alpha)=\sum_{k=0}^{n}\alpha^k\biggl(\frac{(2n-k)!}{k!\,(n-k)!^2}\biggr)^2 .
\end{equation}
The second expression follows from \eqref{eq:apery} with $k=n-j$, since
$\binom nj\binom{n+j}{j}=\frac{(n+j)!}{j!^2(n-j)!}$. Let
$\theta_1,\theta_2,\theta_3$ be independent random variables, uniformly
distributed on $[0,2\pi]$, and let $\E$ denote expectation. Put
\[
Y=4\cos\theta_1\cos\theta_2,\qquad L=Y+b\cos\theta_3 .
\]

\begin{proposition}\label{prop:spectral}
For complex $\tau$ with $|\tau|<1/(2B)$,
\begin{equation}\label{eq:spectral}
\sum_{n\ge0}D_n(\alpha)\,\tau^{2n}=\E\Bigl[\frac{1}{1-\tau L+\alpha\tau^2}\Bigr].
\end{equation}
\end{proposition}

\begin{proof}
Put $\beta=\sqrt\alpha\,\tau$. Then
$1-\tau L+\alpha\tau^2=(1-\beta e^{i\theta_3})(1-\beta e^{-i\theta_3})-\tau Y$,
and
\[
\frac{1}{1-\tau L+\alpha\tau^2}
=\sum_{j\ge0}(\tau Y)^j\sum_{k,l\ge0}\binom{j+k}{j}\binom{j+l}{j}\beta^{k+l}e^{i(k-l)\theta_3}.
\]
Here $|\tau|<1/8$ and $|\beta|<1/8$, and $|Y|\le4$, so the triple series is
dominated by $\sum_j(4|\tau|)^j(1-|\beta|)^{-2j-2}<\infty$, uniformly in the
angles. Taking expectations termwise, $\E[e^{i(k-l)\theta_3}]=\delta_{kl}$,
and $\E[Y^j]=\E[(2\cos\theta_1)^j]^2$ equals $\binom{2i}{i}^2$ if $j=2i$ and
vanishes for odd $j$. Hence
\[
\E\Bigl[\frac{1}{1-\tau L+\alpha\tau^2}\Bigr]
=\sum_{i,k\ge0}\binom{2i}{i}^2\binom{2i+k}{k}^2\alpha^k\tau^{2i+2k}.
\]
The coefficient of $\tau^{2n}$ is
\[
\sum_{k=0}^n\alpha^k\binom{2n-2k}{n-k}^2\binom{2n-k}{k}^2
=\sum_{k=0}^n\alpha^k\biggl(\frac{(2n-k)!}{k!\,(n-k)!^2}\biggr)^2,
\]
which is $D_n(\alpha)$ by \eqref{eq:D}.
\end{proof}

\begin{remark}\label{rem:torus}
Proposition~\ref{prop:spectral} is a torus average of Straub's rational
function \cite[Theorem~3.1]{St}. Substituting $u=\varrho e^{i\theta_1}$,
$v=\varrho e^{-i\theta_1}$, $w=\sigma e^{i\theta_2}$, $z=\sigma e^{-i\theta_2}$
in $1/\bigl((1-u-v)(1-w-z)-\alpha uvwz\bigr)$ and averaging over
$\theta_1,\theta_2$ keeps exactly the monomials with equal exponents within
each pair, and the diagonal in $(\varrho,\sigma)$ of the result is again an
angular average.
\end{remark}

\begin{lemma}\label{lem:f}
The random variable $Y$ has the density $f$ of \eqref{eq:f}. In particular,
$f$ is even, $f>0$ on $(-4,4)$ and $\int f=1$.
\end{lemma}

\begin{proof}
The variables $\cos\theta_i$ are independent with the arcsine density
$\pi^{-1}(1-y^2)^{-1/2}$ on $(-1,1)$. For $0<a<1$ the density of
$\cos\theta_1\cos\theta_2$ at $\pm a$ is
\[
2\int_a^1\frac1r\cdot\frac{1}{\pi\sqrt{1-r^2}}\cdot\frac{1}{\pi\sqrt{1-a^2/r^2}}\dd r
=\frac{2}{\pi^2}\int_a^1\frac{\dd r}{\sqrt{(1-r^2)(r^2-a^2)}} .
\]
The substitution $r^2=1-(1-a^2)\sin^2\phi$ gives
$\dd r/\sqrt{(1-r^2)(r^2-a^2)}=-\dd\phi/r$, so the last integral equals
$\Kc(1-a^2)$. Since $Y=4\cos\theta_1\cos\theta_2$, its density at $c$ is
$\frac14\cdot\frac2{\pi^2}\Kc(1-c^2/16)=f(c)$ for $0<|c|<4$.
\end{proof}

Let $\omega_b(u)=\pi^{-1}(b^2-u^2)^{-1/2}$ for $|u|<b$ and $\omega_b(u)=0$
otherwise; it is the density of $b\cos\theta_3$. Thus $L$ has the density
$h=f*\omega_b$, which is \eqref{eq:h}. Let
\begin{equation}\label{eq:g}
g(w)=\int_{-B}^{B}\frac{h(v)}{w-v}\dd v\qquad(w\in\C\setminus[-B,B])
\end{equation}
be its Cauchy transform.

\begin{lemma}\label{lem:h}
The density $h$ and its Cauchy transform $g$ have the following properties.
\begin{enumerate}[(a)]
\item $h$ is even and nonnegative, vanishes outside $(-B,B)$, is positive on
$(-B,B)$, and $\int h=1$.
\item The function $h$ is H\"older continuous on $\R$ with exponent $1/8$; in
particular $h(\pm B)=0$.
\item $g$ is holomorphic on $\C\setminus[-B,B]$, bounded there, and
$|g(w)|\le2/|w|$ for $|w|\ge2B$. For real $\xi$ the limits
$g(\xi\pm i0)=\lim g(w)$, as $w\to \xi$ within $\Pip$ or $\Pim$ respectively,
exist, depend continuously on $\xi$,
and satisfy
\begin{equation}\label{eq:plemelj}
\re g(\xi+i0)=\re g(\xi-i0),\qquad \im g(\xi\pm i0)=\mp\pi h(\xi).
\end{equation}
\end{enumerate}
\end{lemma}

\begin{proof}
(a) All statements except positivity follow from Lemma~\ref{lem:f}. If
$|w|<B$, the set of $u\in(-b,b)$ with $|w-u|<4$ is a nonempty open interval,
on which $f(w-u)\omega_b(u)>0$.

(b) The function $f$ is bounded by $C(1+|\log|c||)$, satisfies
$|f'(c)|\le C(1+|c|^{-1})$ on $(-4,4)\setminus\{0\}$, by Lemma~\ref{lem:log},
and has jumps at $\pm4$. Let $0<|\delta|\le1$. On the set of points within
$2|\delta|$ of $\{0,\pm4\}$, $|f(c+\delta)-f(c)|\le
C(1+|\log|c||+|\log|c+\delta||)$, which contributes
$O(|\delta|(1+|\log|\delta||)^4)$ to $\|f(\cdot+\delta)-f\|_4^4$. On the
complement in $[-5,5]$, the mean value theorem gives
$|f(c+\delta)-f(c)|\le C|\delta|(1+|c|^{-1})$, which contributes
$O(|\delta|^4\cdot|\delta|^{-3})=O(|\delta|)$. Thus
$\|f(\cdot+\delta)-f\|_4\le C|\delta|^{1/4}(1+|\log|\delta||)$. Since
$\omega_b\in L^{4/3}$, H\"older's inequality yields
$|h(w+\delta)-h(w)|\le\|f(\cdot+\delta)-f\|_4\|\omega_b\|_{4/3}\le
C'|\delta|^{1/8}$.

(c) Holomorphy and the bound for $|w|\ge2B$ are clear. Let $w=\xi+iy$ with
$0<|y|\le1$, and split
\[
g(w)=\int_{|v-\xi |<1}\frac{h(v)-h(\xi)}{w-v}\dd v+h(\xi)\int_{-1}^{1}\frac{\dd s}{s+iy}
+\int_{|v-\xi |\ge1}\frac{h(v)}{w-v}\dd v .
\]
By (b) the first integrand is bounded by $C|v-\xi |^{-7/8}$ and converges
pointwise as $y\to0$. The middle integral equals $-2i\arctan(1/y)$, which
tends to $\mp i\pi$ as $y\to0^\pm$. The last integral is bounded by $1$, and $|g(w)|\le1$ when $|\im w|\ge1$.
Dominated convergence, which is uniform for $\xi$ in compact sets, gives the
boundedness of $g$, the existence and
continuity of the limits, and \eqref{eq:plemelj}, because the limits of the
first and the last term are real.
\end{proof}

\section{Continuation of the spectral density}\label{sec:continuation}

From now on $0<\alpha\le1$, so that $0<b\le2$. For $u\in\C\setminus\bigl((-\infty,-b]\cup[b,\infty)\bigr)$ let
\[
S_b(u)=\sqrt{b-u}\,\sqrt{b+u}
\]
with principal square roots. This is the holomorphic function that equals
$\sqrt{b^2-u^2}>0$ on $(-b,b)$; thus $\omega_b(u)=1/(\pi S_b(u))$ there. For
$w\in\Pim$ let $[w,b]$ denote the oriented segment from $w$ to $b$, and define
\begin{align}
T_-(w)&=\int_{-b}^{b}P(w-u)\,\omega_b(u)\dd u,\label{eq:Tminus}\\
C_-(w)&=\frac1\pi\int_{[w,b]}\Kc\Bigl(\frac{(w-u)^2}{16}\Bigr)\frac{\dd u}{S_b(u)},\label{eq:Cminus}\\
H(w)&=T_-(w)-\frac{i}{\pi^2}\,C_-(w).\label{eq:H}
\end{align}

\begin{lemma}\label{lem:holo}
The functions $T_-$, $C_-$ and $H$ are holomorphic on $\Pim$. For $w\in\Pim$,
\begin{equation}\label{eq:Cparam}
C_-(w)=\frac{\sqrt{b-w}}{\pi}\int_0^1\Kc\Bigl(\frac{(b-w)^2v^2}{16}\Bigr)
\frac{\dd v}{\sqrt{1-v}\,\sqrt{b+w+(b-w)v}} .
\end{equation}
\end{lemma}

\begin{proof}
For real $u$ we have $w-u\in\Pim$, and on compact subsets of $\Pim$ the
integrand of \eqref{eq:Tminus} is holomorphic in $w$ and bounded by a constant
times $\omega_b(u)$, by Lemma~\ref{lem:P}(c). Hence $T_-$ is holomorphic.

Parametrize $[w,b]$ by $u=w+(b-w)v$, $0\le v\le1$. Then $b-u=(b-w)(1-v)$, and
$b-w\in\Pip$, so $\sqrt{b-u}=\sqrt{b-w}\sqrt{1-v}$. Moreover
$\im(b+u)=(1-v)\im w<0$ for $v<1$, while $b+u=2b$ for $v=1$; hence the path
avoids the cuts of $S_b$. Since $w-u=-(b-w)v$ is not real for $0<v\le1$, its
square does not lie in $[0,\infty)$, so the argument of $\Kc$ stays in
$\C\setminus[1,\infty)$. This gives \eqref{eq:Cparam}. On compact subsets of
$\Pim$ the integrand of \eqref{eq:Cparam} is holomorphic in $w$ and bounded by
a constant times $(1-v)^{-1/2}$, so $C_-$ is holomorphic.
\end{proof}

\begin{proposition}[The central interval]\label{prop:central}
$H$ extends continuously to $\Pim\cup(-b,b)$, with $H(\xi)=h(\xi)$ for
$-b<\xi<b$. Consequently $h$ is real-analytic on $(-b,b)$, $H$ extends to a
holomorphic function on $\Pim\cup(-b,b)\cup\Pip$, again denoted $H$, and
\begin{equation}\label{eq:Hsym}
H(w)=\overline{H(-\overline w)}\qquad(w\in\Pim).
\end{equation}
In particular $H(-iy)$ is real for every $y>0$.
\end{proposition}

\begin{proof}
Fix $\xi\in(-b,b)$. For $u\in(-b,b)$ the difference $\xi-u$ lies in
$(\xi-b,\xi+b)$, and since $b\le2$,
\begin{equation}\label{eq:inside}
-4\le-2b<\xi-b<\xi+b<2b\le4 .
\end{equation}
Let $w\to \xi$ within $\Pim$. By Lemma~\ref{lem:P}(a),(b), for $u\ne \xi$,
\[
P(w-u)\longrightarrow
\begin{cases}
f(\xi-u), & u<\xi,\\
f(u-\xi)+\dfrac{i}{\pi^2}\Kc\Bigl(\dfrac{(u-\xi)^2}{16}\Bigr), & u>\xi,
\end{cases}
\]
where we used that $f$ is even and that $u-\xi\in(0,4)$ for $u>\xi$, by
\eqref{eq:inside}. By Lemma~\ref{lem:P}(c), and because $w-u$ stays away from
$-4$, we have $|P(w-u)|\le C(1+|\log|\re w-u||)$. These bounds, multiplied by
$\omega_b$, are uniformly integrable as $\re w$ varies near $\xi$: the
logarithms are bounded in $L^4$ and $\omega_b\in L^{4/3}$. Vitali's theorem
gives
\[
T_-(w)\longrightarrow h(\xi)+\frac{i}{\pi^2}\int_\xi ^b\Kc\Bigl(\frac{(u-\xi)^2}{16}\Bigr)\omega_b(u)\dd u .
\]
In \eqref{eq:Cparam}, as $w\to \xi$, we have $b-w\to b-\xi\in(0,4)$ and
$b+w+(b-w)v\to b+\xi+(b-\xi)v>0$, and the arguments of $\Kc$ converge to
$(b-\xi)^2v^2/16\in[0,1)$, uniformly in $v$. By dominated convergence and the
substitution $u=\xi+(b-\xi)v$,
\[
C_-(w)\longrightarrow\int_\xi ^b\Kc\Bigl(\frac{(u-\xi)^2}{16}\Bigr)\omega_b(u)\dd u .
\]
Hence $H(w)\to h(\xi)$. The convergence is locally uniform in $\xi$, and $h$ is
continuous, so $H$ is continuous on $\Pim\cup(-b,b)$ and real on $(-b,b)$.
The Schwarz reflection principle provides the holomorphic extension, and $h$
is real-analytic on $(-b,b)$ as the restriction of $H$.

The function $w\mapsto\overline{H(-\overline w)}$ is holomorphic on $\Pim$
and has the boundary values $h(-\xi)=h(\xi)$ on $(-b,b)$, by
Lemma~\ref{lem:h}(a). Its difference with $H$ is holomorphic on $\Pim$ with
boundary values $0$ on $(-b,b)$. By reflection, it extends holomorphically
across $(-b,b)$ and vanishes there, hence vanishes identically. This proves
\eqref{eq:Hsym}; for $w=-iy$ we have $-\overline w=w$.
\end{proof}

\begin{remark}\label{rem:xge1}
The inequality $b\le2$, that is $\alpha\le1$ or $x\ge1$, is essential for
Proposition~\ref{prop:central}. By \eqref{eq:inside}, every difference
$\xi-u$ with $\xi,u\in(-b,b)$ lies in $(-4,4)$, so that $f(\xi-u)$ is given
by the single analytic expression \eqref{eq:f}, and the
continuation $P$ of Lemma~\ref{lem:P} applies throughout. If $b>2$, the
differences $\xi-u$ with $\xi,u$ in the central interval can cross $\pm4$,
where $f$ jumps to $0$, and the argument breaks down. The same assumption is
also used in Propositions~\ref{prop:beyond} and~\ref{prop:bounds}, in the
identity $T(b)=h(b)$ and in the lower bounds $|w-u+4|\ge2$ and
$\re(1+dv/4)\ge\frac12$, and, for convenience, in
Proposition~\ref{prop:spectral}. We make no assertion for $x<1$.
\end{remark}

To the right of $b$ we use boundary values from below. Let
\[
Q=\{w\in\C:\ \re w\ge0,\ \im w\le0\}.
\]
For $\xi>b$ we write $H(\xi-i0)=\lim H(w)$ as $w\to \xi$ within $\Pim$, and for
$m\ge0$ we write $\Kc(m-i0)$ for the boundary value from below when $m>1$ and
for $\Kc(m)$ when $m<1$.

\begin{proposition}[Beyond the central interval]\label{prop:beyond}
$H$ extends continuously to $Q$. For $\xi\ge b$,
\begin{equation}\label{eq:Hbeyond}
H(\xi-i0)=T(\xi)+\frac1{\pi^3}\int_b^\xi\Kc\Bigl(\frac{(\xi-u)^2}{16}-i0\Bigr)\frac{\dd u}{\sqrt{u^2-b^2}},
\end{equation}
with $T$ as in \eqref{eq:T}. Consequently $\re H(\xi-i0)=T(\xi)+J(\xi)$ with $J$ as
in \eqref{eq:J}, and for $\xi>b$
\begin{equation}\label{eq:positive}
\begin{gathered}
0\le h(\xi)\le T(\xi),\\
2\re H(\xi-i0)-h(\xi)=2T(\xi)-h(\xi)+2J(\xi)\ge T(\xi)>0 .
\end{gathered}
\end{equation}
Moreover $H(b-i0)=T(b)=h(b)$.
\end{proposition}

\begin{proof}
On $Q\cap\{\re w<b\}$ near the real axis, continuity follows from
Proposition~\ref{prop:central}. Let $\xi\ge b$ and $w\to \xi$ within $\Pim$. For
$u\in(-b,b)$ we have $\xi-u>0$, so by Lemma~\ref{lem:P}(a)
$P(w-u)\to\frac1{2\pi^2}\Kc(1-(\xi-u)^2/16)$. The same uniform integrability
as before gives $T_-(w)\to T(\xi)$.

Consider first the endpoint $\xi=b$. For $u\in(-b,b)$ we have
$0<b-u<2b\le4$, so $f(b-u)=\frac1{2\pi^2}\Kc(1-(b-u)^2/16)$ and hence
$T(b)=h(b)$. Let $w\in\Pim$ with $|w-b|\le\min(b/2,2)$. In
\eqref{eq:Cparam} the arguments of $\Kc$ satisfy
$|(b-w)^2v^2/16|\le\frac14$, where $\Kc$ is bounded, and
$|b+w+(b-w)v-2b|=|w-b|(1-v)\le b/2$, so that $|b+w+(b-w)v|\ge b$.
Since $\int_0^1(1-v)^{-1/2}\dd v=2$, we obtain
$|C_-(w)|\le C|b-w|^{1/2}$, and the same bound holds for the boundary values
on $(b,b+\min(b/2,2)]$. Hence $C_-(w)\to0$ and $H(w)\to T(b)=h(b)$ as $w\to b$
within $Q$. No sign of $\im(b-w)^2$ is used here; indeed this sign is not
constant near $b$, since $w=b-\varepsilon-i\varepsilon$ gives
$(b-w)^2=2i\varepsilon^2$.

Now let $\xi>b$. In \eqref{eq:Cparam}, $b-w$ tends to the negative number
$b-\xi$ from $\Pip$,
so $\sqrt{b-w}\to i\sqrt{\xi-b}$. Also $b+w+(b-w)v\to2b+(\xi-b)(1-v)>0$. Finally
$\im(b-w)^2=2\re(b-w)\im(b-w)<0$ near $\xi$, so the arguments of $\Kc$ approach
$(\xi-b)^2v^2/16$ from the lower half-plane. By Lemma~\ref{lem:bounds} the
integrands are bounded by a constant times
$(1+|\log|1-\lambda v||)(1-v)^{-1/2}$ with $\lambda=|b-w|/4$, and this family
is uniformly integrable for $\lambda$ in compact sets.
Hence
\[
C_-(w)\longrightarrow\frac{i\sqrt{\xi-b}}{\pi}\int_0^1\Kc\Bigl(\frac{(\xi-b)^2v^2}{16}-i0\Bigr)
\frac{\dd v}{\sqrt{1-v}\sqrt{2b+(\xi-b)(1-v)}} ,
\]
and the substitution $u=b+(\xi-b)(1-v)$ turns the right-hand side into
\[
\frac{i}{\pi}\int_b^\xi\Kc\Bigl(\frac{(\xi-u)^2}{16}-i0\Bigr)\frac{\dd u}{\sqrt{u^2-b^2}} .
\]
Multiplying by $-i/\pi^2$ gives \eqref{eq:Hbeyond}. The convergence is
locally uniform in $\xi>b$; together with the endpoint case and
Proposition~\ref{prop:central}, this proves continuity on $Q$.

By Lemma~\ref{lem:cut}, $\re\Kc(m-i0)=\kappa(m)$ for $m\ge0$, $m\ne1$; hence
$\re H(\xi-i0)=T(\xi)+J(\xi)$, and $J\ge0$ because $\kappa>0$. For $\xi>b$ and
$u\in(-b,b)$ we have $\xi-u>0$, and $f(\xi-u)=\frac1{2\pi^2}\Kc(1-(\xi-u)^2/16)$ if
$\xi-u<4$, while $f(\xi-u)=0$ otherwise. Since $\Kc$ is positive on $(-\infty,1)$,
this gives $0\le h(\xi)\le T(\xi)$, and \eqref{eq:positive} follows, with
$T(\xi)>0$. For $\xi=b$ the integral in \eqref{eq:Hbeyond} is empty, and
\eqref{eq:Hbeyond} reduces to $H(b-i0)=T(b)=h(b)$, proved above.
\end{proof}

\begin{proposition}[Bounds]\label{prop:bounds}
For every $R>0$, $H$ is bounded on $Q\cap\{|w|\le R\}$. There are constants
$C$ and $R_0$, depending on $\alpha$, such that
\begin{equation}\label{eq:Hbound}
|H(w)|\le C\,\frac{\log^2|w|}{|w|}\qquad(w\in Q,\ |w|\ge R_0).
\end{equation}
\end{proposition}

\begin{proof}
By continuity it suffices to consider $w\in Q\cap\Pim$. For $u\in(-b,b)$ we
have $\re(w-u)>-b\ge-2$, hence $|w-u+4|\ge2$. If $|w|\le R$,
Lemma~\ref{lem:P}(c) gives $|P(w-u)|\le C_R(1+|\log|\re w-u||)$, and
$\sup_{t\in\R}\int(1+|\log|t-u||)\omega_b(u)\dd u<\infty$; so $T_-$ is
bounded. If $|w|\ge16$, then $|w-u|\ge|w|/2\ge8$ and
$|P(w-u)|\le C\log|w|/|w|$, so $T_-(w)=O(\log|w|/|w|)$.

For $C_-$ write $d=w-b$ and $D=|d|$, so that $\re d\ge-b$. For $0\le v\le1$
the number $\zeta=d(1-v)$ satisfies $\re(2b+\zeta)\ge b$ and
$|2b+\zeta|\ge|\zeta|-2b$; distinguishing $|\zeta|\le3b$ and $|\zeta|>3b$ one
obtains $|2b+\zeta|\ge c_b(1+|\zeta|)$ with $c_b=b/(1+3b)$. Moreover the
argument $m=d^2v^2/16$ of $\Kc$ in \eqref{eq:Cparam} satisfies
$1-m=(1-dv/4)(1+dv/4)$ with $\re(1+dv/4)\ge1-b/4\ge\frac12$ and
$|1-dv/4|\ge|1-Dv/4|$. Therefore \eqref{eq:Cparam} gives
\begin{equation}\label{eq:Cest}
|C_-(w)|\le\frac{1}{\pi\sqrt{c_b}}\int_0^1|\Kc(m)|\,
\frac{\sqrt D\,\dd v}{\sqrt{1-v}\,\sqrt{1+D(1-v)}} .
\end{equation}

If $D\le R$, Lemma~\ref{lem:bounds} and the lower bound for $|1-m|$ give
$|\Kc(m)|\le C_R(1+|\log|1-Dv/4||)$, and
$\sup_{0\le\lambda\le R}\int_0^1(1+|\log|1-\lambda v||)(1-v)^{-1/2}\dd v<\infty$
by H\"older's inequality. Hence $C_-$ is bounded on bounded parts of $Q$.

Let $D\ge16$ and split $[0,1]$ into $I_1=[0,8/D]$, $I_2=[8/D,1-1/D]$ and
$I_3=[1-1/D,1]$.
\begin{itemize}
\item On $I_1$ we have $1-v\ge\frac12$, so the weight in \eqref{eq:Cest} is at
most $2$, and $|m|\le4$, so $|\Kc(m)|\le C(1+|\log|1-Dv/4||)$. The
contribution is at most
$C\int_0^{8/D}(1+|\log|1-Dv/4||)\dd v=O(1/D)$. A possible crossing of
$Dv=4$ contributes only an integrable logarithm.
\item On $I_2$ we have $|m|^{1/2}=Dv/4\ge2$, so
$|\Kc(m)|\le C\log D/(Dv)$, and the weight is at most $(1-v)^{-1}$. The
contribution is at most
$C\frac{\log D}{D}\int_{8/D}^{1-1/D}\frac{\dd v}{v(1-v)}=O\bigl(\frac{\log^2D}{D}\bigr)$.
\item On $I_3$ we have $|m|\ge D^2/64\ge4$, so $|\Kc(m)|\le C\log D/D$, and
the weight is at most $\sqrt D(1-v)^{-1/2}$, whose integral over $I_3$ equals
$2$. The contribution is $O(\log D/D)$.
\end{itemize}
Since $D\ge|w|-b$, we obtain $C_-(w)=O(\log^2|w|/|w|)$, and \eqref{eq:Hbound}
follows from \eqref{eq:H}.
\end{proof}

\section{The Stieltjes transform and the proof of Theorem~\ref{thm:main}}\label{sec:proof}

Let $0<\alpha\le1$, and let $S>\alpha$ be determined by
$\sqrt S+\alpha/\sqrt S=B$, that is,
\begin{equation}\label{eq:S}
S=\biggl(\frac{B+\sqrt{B^2-4\alpha}}{2}\biggr)^2 .
\end{equation}
For $s\in\C\setminus(-\infty,0]$ let $q=q(s)=\sqrt s$ be the principal square
root, and put
\[
r=r(s)=q+\frac{\alpha}{q} .
\]

\begin{lemma}\label{lem:geometry}
The map $s\mapsto r(s)$ has the following properties.
\begin{enumerate}[(a)]
\item Let $s\in\Pip$. Then $\re r>0$ and
$\operatorname{sign}\im r=\operatorname{sign}(|s|-\alpha)$. If $|s|=\alpha$,
then $r\in(0,b)$. If $|s|<\alpha$, then $r$ lies in the interior of $Q$.
\item $r$ maps $(0,\alpha]$ decreasingly onto $[b,\infty)$, $[\alpha,S]$
increasingly onto $[b,B]$, and $(S,\infty)$ onto $(B,\infty)$.
\item As $s\to-a$ from $\Pip$, with $a>0$, we have $q\to i\sqrt a$ and
$r\to i(\sqrt a-\alpha/\sqrt a)$.
\item If $s\in\Pip$ and $|s|\le\alpha/4$, then $|r|\ge3\alpha/(4|q|)$ and
$|r|\,|q|\ge3\alpha/4$.
\end{enumerate}
\end{lemma}

\begin{proof}
For $s\in\Pip$ the root $q$ lies in the open first quadrant, and
$\re r=\re q\,(1+\alpha/|q|^2)>0$, $\im r=\im q\,(1-\alpha/|q|^2)$. If
$|q|=\sqrt\alpha$, then $q=\sqrt\alpha e^{i\phi}$ with $0<\phi<\pi/2$ and
$r=2\sqrt\alpha\cos\phi\in(0,b)$. Parts (b) and (c) are elementary. For (d),
$|q|^2\le\alpha/4$ gives $|q|\le\alpha/(4|q|)$, so
$|r|\ge\alpha/|q|-|q|\ge3\alpha/(4|q|)$; and $|r||q|=|s+\alpha|\ge3\alpha/4$.
\end{proof}

Define $\widehat G$ on $\Pip\cup\Pim$ by
\[
\widehat G(w)=g(w)\quad(w\in\Pip),\qquad
\widehat G(w)=g(w)-2\pi i\,H(w)\quad(w\in\Pim).
\]

\begin{lemma}\label{lem:Ghat}
$\widehat G$ extends holomorphically across $(-b,b)$, and continuously to the
point $b$ from both $\Pip$ and $Q$.
\end{lemma}

\begin{proof}
For $\xi\in(-b,b)$, Lemma~\ref{lem:h}(c) and Proposition~\ref{prop:central} give
the limits $g(\xi+i0)$ from above and $g(\xi-i0)-2\pi ih(\xi)$ from below. By
\eqref{eq:plemelj} these coincide, and they depend continuously on $\xi$.
Hence $\widehat G$ is continuous across $(-b,b)$, and holomorphic there by
Morera's theorem. At $b$, the limit from $\Pip$ is $g(b+i0)$ and the limit
from $Q$ is $g(b-i0)-2\pi iH(b-i0)=g(b-i0)-2\pi ih(b)$, by
Proposition~\ref{prop:beyond}; again these coincide.
\end{proof}

For $s\in\Pip$ put
\begin{equation}\label{eq:M}
M(s)=\frac{\widehat G(r(s))}{q(s)}=
\begin{cases}
g(r)/q, & |s|>\alpha,\\
\bigl(g(r)-2\pi i H(r)\bigr)/q, & |s|<\alpha .
\end{cases}
\end{equation}
By Lemmas~\ref{lem:geometry}(a) and~\ref{lem:Ghat}, $M$ is holomorphic on
$\Pip$; across the semicircle $|s|=\alpha$ the two lines of \eqref{eq:M}
are continuations of each other.

\begin{lemma}\label{lem:germ}
The series $M_\infty(s)=\sum_{n\ge0}D_n(\alpha)s^{-n-1}$ converges for
$|s|>4B^2$, and $M=M_\infty$ on $\Pip\cap\{|s|>4B^2\}$.
\end{lemma}

\begin{proof}
Let $|s|>4B^2$ and $\tau=1/q$, so that $|\tau|<1/(2B)$; then
Proposition~\ref{prop:spectral} gives the convergence. For $v\in[-B,B]$,
\[
\frac{1}{1-\tau v+\alpha\tau^2}=\frac{q^2}{q^2-vq+\alpha}=\frac{q}{r-v} .
\]
Integrating against $h$ and using \eqref{eq:spectral}, we get
$\sum_nD_n(\alpha)s^{-n}=q\,g(r)$. Dividing by $s=q^2$, and noting
$|s|>\alpha$, gives the claim.
\end{proof}

Define $\rho_\alpha$ on $(0,S)$ by
\begin{equation}\label{eq:rho}
\rho_\alpha(t)=
\begin{cases}
\dfrac{h(r(t))}{\sqrt t}, & \alpha\le t<S,\\[2ex]
\dfrac{2T(r(t))-h(r(t))+2J(r(t))}{\sqrt t}, & 0<t<\alpha .
\end{cases}
\end{equation}

\begin{proposition}\label{prop:M}
$M$ extends to a holomorphic function on $\C\setminus[0,S]$ with
$M(\overline s)=\overline{M(s)}$ and $M=M_\infty$ for $|s|>4B^2$. Moreover:
\begin{enumerate}[(a)]
\item every $s_0\in(0,S]$ has a neighbourhood $N$ such that $M$ is bounded on
$N\setminus[0,S]$, and there is $C$ with $|M(s)|\le C(1+\log^2|s|)$ for
$0<|s|\le\alpha/4$, $s\notin[0,S]$;
\item $M(s)=1/s+O(|s|^{-2})$ as $s\to\infty$;
\item for $t\in(0,S)$ the limits $M(t\pm i0)$ exist, are continuous in $t$,
satisfy $M(t-i0)=\overline{M(t+i0)}$, and
$-\frac1\pi\im M(t+i0)=\rho_\alpha(t)$.
\end{enumerate}
\end{proposition}

\begin{proof}
We first show that $M$ extends continuously from $\Pip$ to
$(-\infty,0)\cup(S,\infty)$ with real values. For $t>S$, $r(t)\in(B,\infty)$
by Lemma~\ref{lem:geometry}(b), where $g$ is holomorphic and real; so
$M(t)=g(r(t))/\sqrt t$ is real. Let $s\to-a$ with $a>0$, and put
$y=\sqrt a-\alpha/\sqrt a$. By Lemma~\ref{lem:geometry}(c),
$M(s)\to\widehat G(iy)/(i\sqrt a)$, where $\widehat G$ is continuous at $iy$
(for $y=0$ by Lemma~\ref{lem:Ghat}). If $y\ne0$, the evenness of $h$ gives
\[
g(iy)=\int\frac{h(v)(-iy-v)}{y^2+v^2}\dd v=-iy\int\frac{h(v)}{y^2+v^2}\dd v\in i\R .
\]
If $y>0$, then $\widehat G(iy)=g(iy)\in i\R$. If $y<0$, then
$\widehat G(iy)=g(iy)-2\pi iH(iy)\in i\R$, because $H(iy)$ is real by
Proposition~\ref{prop:central}. If $y=0$, $\widehat G(0)\in i\R$ by
continuity. In all cases $M(-a)$ is real. By the Schwarz reflection principle,
$M$ extends holomorphically to $\C\setminus[0,S]$ with
$M(\overline s)=\overline{M(s)}$. Lemma~\ref{lem:germ}, the reality of the
coefficients $D_n(\alpha)$ and continuity give $M=M_\infty$ for $|s|>4B^2$,
and (b) follows because $D_0(\alpha)=1$.

(a) Let $s_0\in(0,S]$ and let $s\in\Pip$ be close to $s_0$. Then $|q|$ is
bounded below. If $|s|\ge\alpha$, then $|\widehat G(r)|$ is bounded by
Lemmas~\ref{lem:h}(c) and~\ref{lem:Ghat}. If $|s|<\alpha$, then $r\in Q$ is
bounded, and $|\widehat G(r)|\le|g(r)|+2\pi|H(r)|$ is bounded by
Lemma~\ref{lem:h}(c) and Proposition~\ref{prop:bounds}. By reflection the same
holds in $\Pim$. Near $0$, let $s\in\Pip$ with $|s|\le\alpha/4$ and so small
that $|r|\ge\max(2B,R_0)$, which is possible by Lemma~\ref{lem:geometry}(d).
Then $|g(r)|\le2/|r|$ and $|H(r)|\le C\log^2|r|/|r|$. By
Lemma~\ref{lem:geometry}(d), $|r||q|\ge3\alpha/4$, and
$|r|\le|q|+\alpha/|q|\le2\alpha/|q|$, so $\log|r|\le C(1+|\log|s||)$. Hence
\[
|M(s)|\le\frac{2+2\pi C\log^2|r|}{|r|\,|q|}\le C'(1+\log^2|s|) .
\]
On the remaining compact part of $\{|s|\le\alpha/4\}$ the bound follows from
the first assertion of (a), and the case $\im s<0$ follows by reflection.

(c) Let $t\in(0,S)$ and $s\to t$ from $\Pip$; then $q\to\sqrt t$. If
$t>\alpha$, then by Lemma~\ref{lem:geometry}(a),(b) the point $r(s)$ tends to
$r(t)\in(b,B)$ from $\Pip$, so $M(t+i0)=g(r(t)+i0)/\sqrt t$, and
\eqref{eq:plemelj} gives $-\frac1\pi\im M(t+i0)=h(r(t))/\sqrt t$. If
$t<\alpha$, then $r(s)$ tends to $r(t)\in(b,\infty)$ within $Q$. Since $H$
is continuous on $Q$ (Proposition~\ref{prop:beyond}),
\[
M(t+i0)=\frac{g(r(t)-i0)-2\pi i\,H(r(t)-i0)}{\sqrt t} .
\]
By \eqref{eq:plemelj}, $\im g(r-i0)=\pi h(r)$, and
$\im(-2\pi iH)=-2\pi\re H$. Therefore
\[
-\frac1\pi\im M(t+i0)=\frac{2\re H(r(t)-i0)-h(r(t))}{\sqrt t},
\]
which equals the second line of \eqref{eq:rho} by
Proposition~\ref{prop:beyond}. If $t=\alpha$, then $r(s)\to b$, and
Lemma~\ref{lem:Ghat} gives $M(\alpha+i0)=g(b+i0)/\sqrt\alpha$. Both lines of
\eqref{eq:rho} tend to $h(b)/\sqrt\alpha$ as $t\to\alpha$, since $T(b)=h(b)$
and $J(b)=0$. Continuity in $t$ follows from the continuity of $g(\cdot\pm
i0)$, of $H$ on $Q$, and of $\widehat G$. The relation
$M(t-i0)=\overline{M(t+i0)}$ is a consequence of $M(\overline
s)=\overline{M(s)}$.
\end{proof}

\begin{theorem}\label{thm:alpha}
Let $0<\alpha\le1$. The function $\rho_\alpha$ of \eqref{eq:rho} is strictly
positive on $(0,S)$, satisfies $\rho_\alpha(t)=O(1+\log^2t)$ as $t\to0^+$, and
is integrable on $(0,S)$. Moreover
\begin{equation}\label{eq:cauchy}
M(s)=\int_0^S\frac{\rho_\alpha(t)}{s-t}\dd t\qquad(s\in\C\setminus[0,S]).
\end{equation}
Consequently $\int_0^S\rho_\alpha(t)\dd t=1$ and
\begin{equation}\label{eq:momentsD}
D_n(\alpha)=\int_0^St^n\rho_\alpha(t)\dd t\qquad\text{for all } n\ge0 .
\end{equation}
\end{theorem}

\begin{proof}
For $\alpha\le t<S$ we have $r(t)\in[b,B)$, where $h>0$ by
Lemma~\ref{lem:h}(a). For $0<t<\alpha$, \eqref{eq:positive} gives
$\rho_\alpha(t)\ge T(r(t))/\sqrt t>0$. Since
$\rho_\alpha(t)=-\frac1\pi\im M(t+i0)$, Proposition~\ref{prop:M}(a) gives
$|\rho_\alpha(t)|\le C(1+\log^2t)$ near $0$, and local boundedness on
$[\varepsilon,S)$; hence $\rho_\alpha$ is integrable.

Fix $s_0\in\C\setminus[0,S]$. For $0<\varepsilon<\min(\alpha/4,1)$ let
$U_\varepsilon=\{z:\operatorname{dist}(z,[0,S])<\varepsilon\}$, and let
$R>|s_0|+S+1$. Cauchy's integral formula on
$\{|z|<R\}\setminus\overline{U_\varepsilon}$, for $\varepsilon$ small, gives
\[
M(s_0)=\frac1{2\pi i}\oint_{|z|=R}\frac{M(z)}{z-s_0}\dd z
-\frac1{2\pi i}\oint_{\partial U_\varepsilon}\frac{M(z)}{z-s_0}\dd z ,
\]
both contours being oriented counterclockwise. The first integral is
$O(1/R)$ by Proposition~\ref{prop:M}(b) and tends to $0$ as $R\to\infty$. The
boundary $\partial U_\varepsilon$ consists of the segments $[0,S]\pm
i\varepsilon$ and of two semicircles of radius $\varepsilon$ centred at $0$
and $S$. By Proposition~\ref{prop:M}(a), the integral over the semicircle at
$S$ is $O(\varepsilon)$, and the integral over the semicircle at $0$ is
$O(\varepsilon\log^2\varepsilon)$; both tend to $0$. On the segments, the
integrands converge pointwise for $t\in(0,S)$ to $M(t\pm i0)/(t-s_0)$, by
Proposition~\ref{prop:M}(c), and they are dominated by $C(1+\log^2t)$, by
Proposition~\ref{prop:M}(a). In the counterclockwise orientation the lower
segment is traversed from $0$ to $S$ and the upper one from $S$ to $0$. By
dominated convergence,
\[
-\frac1{2\pi i}\oint_{\partial U_\varepsilon}\frac{M(z)}{z-s_0}\dd z
\longrightarrow-\frac1{2\pi i}\int_0^S\frac{M(t-i0)-M(t+i0)}{t-s_0}\dd t .
\]
By Proposition~\ref{prop:M}(c),
$M(t-i0)-M(t+i0)=-2i\im M(t+i0)=2\pi i\rho_\alpha(t)$, and the right-hand
side equals $\int_0^S\rho_\alpha(t)(s_0-t)^{-1}\dd t$. This proves
\eqref{eq:cauchy}. In particular, no atom, pole or additional holomorphic
term is present.

Multiplying \eqref{eq:cauchy} by $s$ and letting $s\to+\infty$, dominated
convergence and Proposition~\ref{prop:M}(b) give $\int_0^S\rho_\alpha=1$. For
$|s|>S$ we may expand $(s-t)^{-1}=\sum_{n\ge0}t^ns^{-n-1}$ uniformly in
$t\in[0,S]$, so $M(s)=\sum_{n\ge0}s^{-n-1}\int_0^St^n\rho_\alpha(t)\dd t$.
Comparing with $M=M_\infty$ for $|s|>4B^2$ and using the uniqueness of
Laurent coefficients gives \eqref{eq:momentsD}.
\end{proof}

\begin{proof}[Proof of Theorem~\ref{thm:main}]
Let $x\ge1$ and $\alpha=1/x\in(0,1]$. Then $b=2\sqrt\alpha=2x^{-1/2}$, in
accordance with the definitions of the introduction. By \eqref{eq:S},
\[
\sqrt{S/\alpha}=\frac{B}{2\sqrt\alpha}+\sqrt{\frac{B^2}{4\alpha}-1}
=(1+2\sqrt x)+\sqrt{(1+2\sqrt x)^2-1}
\]
and the right-hand side equals $1+2\sqrt x+2\sqrt{x+\sqrt x}$,
so $S/\alpha=c_x$. Let $\mu_x$ be the image of $\rho_\alpha(t)\dd t$ under
$t\mapsto t/\alpha$. Its density is $\alpha\rho_\alpha(\alpha t)$ on
$(0,c_x)$. Since $r(\alpha t)=\sqrt{\alpha t}+\sqrt{\alpha/t}=r_x(t)$ and
$\alpha/\sqrt{\alpha t}=1/\sqrt{xt}$, and since $\alpha t<\alpha$ exactly when
$t<1$, this density is the function $\varphi_x$ of \eqref{eq:density}. By
\eqref{eq:D} and Theorem~\ref{thm:alpha},
\[
A_n(x)=x^nD_n(1/x)=\alpha^{-n}\int_0^St^n\rho_\alpha(t)\dd t
=\int_0^{c_x}t^n\varphi_x(t)\dd t
\]
for all $n\ge0$. By Theorem~\ref{thm:alpha}, $\mu_x$ is a probability
measure, $\varphi_x$ is strictly positive on $(0,c_x)$, so that the support of
$\mu_x$ is $[0,c_x]$, and $\varphi_x(t)=O(\log^2t)$ as $t\to0^+$.

For uniqueness, let $\nu$ be a positive Borel measure on $\R$ with
$\int t^n\dd\nu=A_n(x)$ for all $n$. For $\lambda>c_x$ we get
$\nu(\{|t|\ge\lambda\})\le\lambda^{-2n}A_{2n}(x)\le(c_x/\lambda)^{2n}\to0$,
so $\nu$ is supported in $[-c_x,c_x]$. There, polynomials are dense in the
continuous functions, and $\nu$ and $\mu_x$ agree on polynomials; hence
$\nu=\mu_x$.
\end{proof}

\section{Further remarks}\label{sec:remarks}

\begin{remark}[Consequences]
For each fixed real $x\ge1$, the Hankel matrix $(A_{i+j}(x))_{i,j\ge0}$ is
totally positive \cite{WZ}, and, by the theorem of Wang and Zhu that
Stieltjes moment sequences are infinitely log-convex \cite{WZ}, the sequence
$(A_n(x))_{n\ge0}$ is infinitely log-convex. The coefficientwise statement
of \cite[Conjecture~2]{So25} does not follow.
\end{remark}

\begin{remark}[Lattice Green functions]
Since $4\cos\theta_1\cos\theta_2=2\cos(\theta_1+\theta_2)+2\cos(\theta_1-\theta_2)$
and $(\theta_1,\theta_2)\mapsto(\theta_1+\theta_2,\theta_1-\theta_2)$ preserves
Haar measure on the torus, for fixed real $x\ge1$ the variable $L$ has the law of
$2(\cos\eta_1+\cos\eta_2+\sqrt\alpha\cos\eta_3)$ with independent uniform
$\eta_i$. Thus $g(w)=\E[(w-L)^{-1}]$ is, up to normalization, the Green
function at the origin of the anisotropic simple cubic lattice with
anisotropy parameter $\sqrt\alpha=x^{-1/2}$, studied by Delves and Joyce
\cite{DJ01,DJ06}. In this language $f$ is the density of states of the square
lattice.
\end{remark}

\section*{Use of AI-assisted tools}
AI-assisted tools were used in the research and preparation of this manuscript.

\begin{thebibliography}{10}

\bibitem{Ap}
R.~Ap\'ery, \emph{Irrationalit\'e de $\zeta(2)$ et $\zeta(3)$},
Ast\'erisque \textbf{61} (1979), 11--13.

\bibitem{DJ01}
R.~T. Delves and G.~S. Joyce, \emph{On the Green function for the anisotropic
simple cubic lattice}, Ann. Physics \textbf{291} (2001), 71--133.

\bibitem{DJ06}
R.~T. Delves and G.~S. Joyce, \emph{Exact product form for the anisotropic
simple cubic lattice Green function}, J. Phys. A \textbf{39} (2006), no.~16,
4119--4145.

\bibitem{DLMF}
NIST Digital Library of Mathematical Functions, F.~W.~J. Olver et al.\
(eds.), \url{https://dlmf.nist.gov/}, accessed October 2026.

\bibitem{Ed}
G.~A. Edgar, \emph{The Ap\'ery numbers as a Stieltjes moment sequence},
preprint (2020), arXiv:2005.10733v2.

\bibitem{Ne}
T.~Neuschel, \emph{Ap\'ery polynomials and the multivariate saddle point
method}, preprint (2013), arXiv:1307.0341.

\bibitem{So20}
A.~D. Sokal, \emph{The Euler and Springer numbers as moment sequences},
Expo. Math. \textbf{38} (2020), 1--26.

\bibitem{So25}
A.~D. Sokal, \emph{Stieltjes moment sequences and coefficientwise
Hankel-total positivity in enumerative combinatorics}, lecture slides,
conference in honour of the 80th birthday of Tony Guttmann, 1 July 2025,
\url{https://clisby.net/guttmann2025/talks/Alan_Sokal.pdf}.

\bibitem{St}
A.~Straub, \emph{Multivariate Ap\'ery numbers and supercongruences of
rational functions}, Algebra Number Theory \textbf{8} (2014), no.~8,
1985--2007.

\bibitem{WZ}
Y.~Wang and B.-X. Zhu, \emph{Log-convex and Stieltjes moment sequences},
Adv. in Appl. Math. \textbf{81} (2016), 115--127.

\end{thebibliography}

\end{document}
